\chapter{Surgical Thrombectomy \& Embolectomy}

\section*{Introduction \& Clinical Indications}

Surgical thrombectomy and embolectomy are critical procedures aimed at the removal of thrombi and emboli from the vascular system, respectively. Thrombectomy refers to the surgical removal of a blood clot that obstructs blood flow, while embolectomy involves the extraction of an embolus, which may be a thrombus or other material that has traveled through the bloodstream and lodged in a vessel. These procedures are essential in the management of acute limb ischemia, myocardial infarction, and pulmonary embolism, among other conditions.

Historically, the development of thrombectomy techniques can be traced back to the early 20th century, with significant advancements in surgical techniques and instrumentation over the decades. The clinical indications for these procedures include:

\begin{itemize}
    \item Acute limb ischemia due to arterial occlusion.
    \item Myocardial infarction with ST-segment elevation.
    \item Pulmonary embolism with hemodynamic instability.
    \item Cerebral embolism leading to acute ischemic stroke.
    \item Venous thrombosis with significant complications.
    \item Complications of vascular grafts or stents.
    \item Trauma-related vascular injuries.
    \item Postoperative thrombosis in high-risk patients.
\end{itemize}

Thrombectomy or embolectomy is selected according to the vascular bed, cause, severity, time course, imaging, and available endovascular or surgical expertise. Acute limb ischemia requires immediate vascular assessment and anticoagulation when appropriate; treatment may include open thrombectomy, catheter-directed thrombolysis, aspiration or mechanical thrombectomy, bypass, or hybrid repair. Myocardial infarction is generally managed with urgent percutaneous coronary intervention rather than open surgical embolectomy.

\section*{Relevant Surgical Anatomy}

A thorough understanding of the relevant surgical anatomy is crucial for the successful execution of thrombectomy and embolectomy procedures. 

The arterial supply to the limbs primarily arises from the subclavian artery, which branches into the axillary, brachial, radial, and ulnar arteries in the upper extremity, and from the aorta, which bifurcates into the common femoral artery, leading to the superficial and deep femoral arteries, popliteal artery, and subsequently the anterior and posterior tibial arteries in the lower extremity.

Venous drainage follows a similar pattern, with deep veins accompanying the arteries and superficial veins draining into the deep system. The venous system is equipped with valves that prevent retrograde flow, which can be compromised in the presence of thrombosis.

Nerve relations are also significant, as the major nerves of the limbs, including the median, ulnar, and radial nerves in the upper extremity, and the sciatic, femoral, and tibial nerves in the lower extremity, must be preserved during surgical interventions.

Lymphatic drainage is primarily through the lymphatic vessels that accompany the venous system, and surgical landmarks such as the inguinal ligament, popliteal fossa, and anatomical landmarks of the arm are critical for safe access during the procedure.

\section*{Preoperative Workup \& Preparation}

A comprehensive preoperative workup is essential to optimize patient outcomes and minimize complications. 

Laboratory tests should include:

\begin{itemize}
    \item Complete blood count (CBC) to assess for anemia or infection.
    \item Coagulation profile (PT, aPTT, INR) to evaluate bleeding risk.
    \item Electrolytes and renal function tests to assess metabolic status.
    \item Type and crossmatch for potential blood transfusion.
\end{itemize}

Imaging studies, such as Doppler ultrasound, CT angiography, or MRI, may be utilized to evaluate the extent of vascular occlusion and to plan the surgical approach.

Cardiovascular optimization is crucial, particularly in patients with underlying cardiovascular disease. This may involve the management of hypertension, diabetes, and other comorbidities.

Informed consent must be obtained, detailing the risks, benefits, and alternatives to the procedure.

Prophylactic antibiotics should be administered to reduce the risk of postoperative infections, and patients should be instructed to remain NPO (nothing by mouth) for a specified period prior to surgery.

Situations requiring individualized risk-benefit assessment include:

\begin{itemize}
    \item Uncontrolled bleeding diathesis.
    \item Severe comorbidity or bleeding risk may alter the procedure, anesthesia, or use of thrombolysis, but does not automatically preclude urgent revascularization.
\end{itemize}

Relative contraindications include:

\begin{itemize}
    \item Advanced age with multiple comorbidities.
    \item Severe peripheral neuropathy.
    \item Infection at the surgical site.
\end{itemize}

\section*{Surgical Technique \& Procedure}

The surgical technique for thrombectomy and embolectomy varies based on the anatomical location and the nature of the occlusion. 

Patient positioning is critical; the patient is typically placed supine on the operating table, with the affected limb positioned for optimal access.

General anesthesia is commonly used, although regional anesthesia may be appropriate in select cases.

The incision is made over the affected vessel, with careful dissection to expose the artery. 

The following steps outline a general approach to the procedure:

\begin{itemize}
    \item Identify the occluded vessel and assess collateral circulation.
    \item Make an incision over the artery, typically along the course of the vessel.
    \item Carefully dissect through the subcutaneous tissue and fascia to expose the artery.
    \item Clamp the artery proximally and distally to control blood flow.
    \item Use a thrombectomy or embolectomy catheter to remove the clot or embolus. This may involve mechanical fragmentation or suction techniques.
    \item Flush the artery with saline or heparinized saline to ensure patency and remove residual debris.
    \item Inspect the vessel for any residual stenosis or injury.
    \item Close the incision in layers, ensuring hemostasis and proper alignment of tissues.
\end{itemize}

Closure involves meticulous suturing of the skin and subcutaneous layers, with attention to hemostasis to minimize the risk of hematoma formation.

\section*{Complications \& Risks}

As with any surgical procedure, thrombectomy and embolectomy carry inherent risks. General surgical complications may include:

\begin{itemize}
    \item Infection at the surgical site.
    \item Hematoma formation.
    \item Wound dehiscence.
    \item Anesthesia-related complications.
\end{itemize}

Procedure-specific risks include:

\begin{itemize}
    \item Re-occlusion of the vessel.
    \item Embolization of debris to distal vessels.
    \item Vascular injury or transection.
    \item Neurological deficits in cases of cerebral embolism.
    \item Compartment syndrome due to reperfusion injury.
    \item Thrombosis of adjacent vessels.
\end{itemize}

Close monitoring during and after the procedure is essential to identify and manage these complications promptly.

\section*{Postoperative Care \& Recovery}

Postoperative care begins in the Post Anesthesia Care Unit (PACU), where patients are monitored for vital signs, neurological status, and signs of complications.

The expected course includes:

\begin{itemize}
    \item Pain management with analgesics as needed.
    \item Monitoring for signs of infection or hematoma.
    \item Assessment of limb perfusion, including capillary refill and pulse checks.
    \item Gradual mobilization as tolerated, with physical therapy initiated as appropriate.
\end{itemize}

Surgical findings and pathology reports should be reviewed to guide further management.

Follow-up care guidelines typically involve:

\begin{itemize}
    \item Regular follow-up visits to assess vascular status.
    \item Doppler studies to evaluate blood flow.
    \item Long-term anticoagulation therapy in cases of thromboembolic disease.
\end{itemize}

\section*{Outcomes \& Prognosis}

Outcomes depend on ischemia severity, cause, tissue viability, time to reperfusion, anatomy, comorbidity, and whether additional reconstruction is required. Early treatment improves the chance of limb salvage, but no universal success percentage applies; reperfusion injury, re-occlusion, bleeding, compartment syndrome, amputation, and recurrent embolism remain possible.

Recurrence risk is influenced by the underlying etiology of thrombosis, with patients requiring ongoing management of risk factors such as hyperlipidemia, hypertension, and diabetes.

Epidemiological studies indicate that early intervention is crucial in improving outcomes, and landmark trials have demonstrated the efficacy of surgical intervention in selected populations.

\section*{Additional Clinical Considerations}
The urgency and treatment are determined by the vascular bed. In acute limb ischemia, the clinician documents the onset, pain, pallor, pulselessness, paresthesia, paralysis, and poikilothermia, while distinguishing an immediately threatened limb from a viable or irreversibly damaged one. Intravenous unfractionated heparin is commonly started when not contraindicated, and urgent vascular consultation, Doppler assessment, and imaging are coordinated without delaying revascularization in a threatened limb. A limb with irreversible muscle or nerve injury may not benefit from revascularization and may require amputation after multidisciplinary assessment.

The cause of an arterial occlusion must be addressed. Embolism from atrial fibrillation, mural thrombus, aneurysm, endocarditis, or a cardiac or aortic source requires investigation and secondary prevention. In-situ thrombosis often reflects peripheral arterial disease, graft failure, stent restenosis, dissection, or hypercoagulability and may require angioplasty, stenting, bypass, or revision in addition to clot removal. Removing the clot without treating the underlying lesion increases the risk of early re-occlusion.

Open catheter embolectomy is not interchangeable with treatment for every thrombosis. Cerebral large-vessel stroke is generally treated with endovascular thrombectomy in an appropriately equipped stroke system, while ST-elevation myocardial infarction is usually managed with urgent percutaneous coronary intervention. Massive or high-risk pulmonary embolism may require catheter-directed therapy, systemic thrombolysis, surgical embolectomy, or mechanical circulatory support according to hemodynamic status and contraindications. Venous thrombosis is usually managed with anticoagulation; thrombectomy is reserved for selected extensive or limb-threatening cases.

Reperfusion can cause edema, hyperkalemia, acidosis, rhabdomyolysis, renal injury, and compartment syndrome. Serial pulse, Doppler, motor, sensory, compartment, urine-output, and laboratory assessments are therefore essential after flow is restored. Bleeding risk is increased by heparin or thrombolysis and must be balanced against re-occlusion. Distal embolization, intimal injury, dissection, pseudoaneurysm, and access-site hematoma may require imaging or further intervention.

Long-term care includes anticoagulation or antiplatelet therapy when indicated, statin treatment, smoking cessation, diabetes and blood-pressure control, atrial-fibrillation management, and surveillance of grafts or stents. The discharge plan identifies the treated vessel, residual disease, medication duration, wound care, and symptoms requiring urgent review. Sudden recurrent pain, pallor, coolness, numbness, weakness, bleeding, chest pain, or breathlessness should be treated as an emergency.

\section*{Documentation and Safety}
The procedure record should identify the occluded vessel, duration and severity of ischemia, imaging, anticoagulation, technique, distal runoff, residual stenosis, pulses or Doppler signals before and after treatment, and any adjunctive repair. Postoperative handover should state the antithrombotic plan, limb or neurologic monitoring, compartment-syndrome risk, and the investigation of the embolic source.

\section*{Practical Follow-Up}
After revascularization, surveillance includes the treated vessel, distal perfusion, bleeding, renal function, and reperfusion injury. The patient needs a secondary-prevention plan for the underlying embolic or atherosclerotic cause. Sudden recurrent ischemic symptoms, bleeding, neurologic change, chest pain, or breathlessness require emergency assessment.

\section*{Safety Summary}
Revascularization is time-sensitive, but the treatment must match the vascular bed and cause. Restoration of flow is followed by surveillance for bleeding, reperfusion injury, re-occlusion, and the underlying embolic or thrombotic disorder.

\section*{Final Safety Note}
Restoring flow is only the first step. Ongoing limb, neurologic, cardiac, or respiratory monitoring and treatment of the underlying embolic or thrombotic cause are essential to prevent re-occlusion and recurrent events.

\section*{Completion Checklist}
Record the treated vessel, ischemia severity, pre- and post-procedure perfusion, anticoagulation plan, reperfusion monitoring, and evaluation for the embolic source. Sudden recurrent ischemia, neurologic change, bleeding, chest pain, or breathlessness is emergent.

\section*{Handoff}
The vascular handoff should state perfusion findings, antithrombotic therapy, compartment monitoring, and the plan for treating the underlying embolic source.

\section*{References}

\begin{enumerate}
    \item Sabiston, D. C., \& Spencer, F. C. (2016). \textit{Surgery of the Chest}. 9th ed. Philadelphia: Elsevier.
    
    \item Schwartz, S. I., \& Shires, G. T. (2019). \textit{Principles of Surgery}. 11th ed. New York: McGraw-Hill.
    
    \item American College of Cardiology/American Heart Association. (2014). \textit{Guidelines for the Management of Patients with ST-Elevation Myocardial Infarction}.
    
    \item Kahn, S. R., et al. (2016). \textit{Venous Thromboembolism: A Clinical Practice Guideline}. Chest, 149(2), 315-352.
\end{enumerate}

\clearpage
